\documentclass[11pt,a4paper]{article}

\usepackage[T1]{fontenc}
\usepackage[utf8]{inputenc}
\usepackage{mathpazo}
\usepackage[margin=2.5cm]{geometry}
\usepackage{microtype}
\usepackage{amsmath,amssymb}
\usepackage{booktabs}
\usepackage{tabularx}
\usepackage{array}
\usepackage{enumitem}
\usepackage{xcolor}
\usepackage[hidelinks]{hyperref}
\usepackage{caption}

\setlist{itemsep=2pt,topsep=4pt}
\captionsetup{font=small,labelfont=bf}
\setlength{\parskip}{4pt}
\setlength{\parindent}{0pt}
\renewcommand{\arraystretch}{1.15}

\newcommand{\code}[1]{\texttt{#1}}
\newcolumntype{L}{>{\raggedright\arraybackslash}X}

\begin{document}

\begin{center}
{\LARGE\bfseries Deliverable 6: Analytical Note}\\[4pt]
{\large Student-to-Project Assignment Decision Support System}\\[10pt]
TECH5 Programming and Application Design, NHH Norwegian School of Economics\\[8pt]
\begin{tabular}{ll}
\textbf{Group:} & 15\\
\textbf{Team B (Assignment Engine):} & Derek Liang, Aayush [surname], Ole Theodor [surname]\\
\textbf{Team A (Visualization):} & [names]\\
\textbf{Date:} & [submission date]
\end{tabular}
\end{center}

\vspace{6pt}
\noindent\textbf{Summary.}
We assign students to projects with an integer linear program (ILP). All
business rules that must never be broken (one project per student, capacities,
fixed assignments, mandatory skills) are hard constraints. Three strategies
change only the objective: preference, skill fit, and a balanced strategy that
first protects the worst-off student. Up to $K$ distinct alternatives are
generated with no-good cuts, and impossible problems are reported as
infeasible rather than silently relaxed. Each alternative is evaluated on
feasibility, efficiency, fairness and capacity, and recommendations are made
per policy priority so that the decision-maker sees the trade-off instead of a
single score.

\section{Problem definition and modelling approach}

\subsection{Business context and decision dilemma}

A programme coordinator must assign a cohort of students to semester-long
business projects, as in the CEMS Business Projects at NHH. Students rank the
projects they prefer and have skills at levels 1--5; projects have a minimum
and maximum number of students, desired skills and some mandatory skills; and
the administration may already have fixed some assignments. The objectives
conflict: giving many students their first choice can leave a project without
the skills it needs, and maximising average satisfaction can leave one student
with a very poor outcome. Rezaeinia, G\'oez and Guajardo~\cite{rezaeinia2022}
study this setting for an NHH master's programme and show that efficiency and
fairness have to be treated as separate criteria. Following that insight, our
system does not hide the conflict behind one score: it produces several
feasible alternatives that each favour a different criterion, and it reports
indicators that make the differences visible.

\subsection{Chosen methodology and justification}

We use binary integer programming, modelled with Pyomo~\cite{pyomo} and solved
with the HiGHS solver~\cite{highs} (\code{appsi\_highs}), the same tools as in
TECH4. The model is a generalised
\emph{assignment problem}~\cite[Ch.~9.3]{ior}: every student is assigned to
exactly one project, but a project receives between $l_p$ and $u_p$ students
instead of exactly one. Additional business rules are added with the standard
binary-variable techniques for logical relationships~\cite[Ch.~12]{ior}.

We chose ILP over the alternatives the brief mentions for four reasons.
\begin{enumerate}
  \item \textbf{Every rule maps to one constraint.} Capacities, fixed
        assignments and mandatory skills are linear constraints, which makes
        the model easy to inspect and explain. A greedy rule-based procedure
        needs repair steps for minimum capacities and mandatory skills, and
        stable matching loses its guarantees once lower quotas are added.
  \item \textbf{Infeasibility is proven.} When no assignment satisfies the
        hard constraints, branch-and-bound~\cite[Ch.~12.6]{ior} proves it. A
        heuristic that fails cannot distinguish ``impossible'' from ``not
        found''.
  \item \textbf{Distinct alternatives are exact.} No-good cuts
        (constraint~\eqref{eq:nogood}) guarantee that every alternative
        differs from all earlier ones in at least one student's project, which
        is exactly the brief's definition of a distinct alternative.
  \item \textbf{The problem is small.} The model has $|S|\cdot|P| + |P|$
        binary variables: 44 for the 10-student demo cohort, which HiGHS solves
        in a few hundredths of a second inside the single-process Flask
        application.
\end{enumerate}

\section{Mathematical formulation}

\subsection{Index sets}
\begin{itemize}
  \item $S$: students in the cohort, $|S| = n$.
  \item $P$: projects, $|P| = m$.
  \item $K_p$: skills required by project $p$ (mandatory and desired);
        $M_p \subseteq K_p$: the mandatory ones.
  \item $F \subseteq S \times P$: fixed administrative assignments;
        $S_F = \{ s : (s,p) \in F \}$: students with a fixed assignment.
  \item $\mathcal{A} = \{A^{(1)}, \dots, A^{(j)}\}$: solutions found earlier
        in the same run, each a set of $n$ student--project pairs.
\end{itemize}

\subsection{Parameters}
\begin{itemize}
  \item $l_p \in \mathbb{Z}_{\ge 0}$, $u_p \in \mathbb{Z}_{\ge 1}$: minimum
        and maximum number of students on project $p$, $l_p \le u_p$.
  \item $r_{sp}$: preference rank of student $s$ for project $p$
        (1 = first choice). An unranked project gets $r_{sp} = m + 1$.
  \item $u_{sp} = (m + 1 - r_{sp})/m \in [0,1]$: preference score, 1 for a
        first choice and 0 for an unranked project.
  \item $\lambda_{sk} \in \{0,\dots,5\}$: level of student $s$ in skill $k$
        (0 if not held); $L_{pk} \in \{1,\dots,5\}$: level project $p$
        requires in skill $k$.
  \item $\sigma_{sp} \in [0,1]$: skill fit of student $s$ on project $p$, the
        average share of each required level that the student reaches:
        \begin{equation}
          \sigma_{sp} =
          \begin{cases}
            \dfrac{1}{|K_p|} \displaystyle\sum_{k \in K_p}
              \min\!\left(1, \dfrac{\lambda_{sk}}{L_{pk}}\right)
              & \text{if } K_p \neq \emptyset,\\[8pt]
            1 & \text{if } K_p = \emptyset.
          \end{cases}
          \label{eq:fit}
        \end{equation}
        A project without requirements is a neutral, perfect fit. A student
        with no relevant skills has fit 0 but can still be assigned.
  \item $\varepsilon = 0.001$: tie-break weight (Section~\ref{sec:obj}).
\end{itemize}

\subsection{Decision variables}
\begin{align*}
  x_{sp} &= \begin{cases} 1 & \text{if student $s$ is assigned to project $p$,}\\
                           0 & \text{otherwise,} \end{cases}
  & \forall s \in S,\ p \in P,\\[4pt]
  y_p &= \begin{cases} 1 & \text{if project $p$ runs (has at least one student),}\\
                        0 & \text{otherwise,} \end{cases}
  & \forall p \in P.
\end{align*}
The balanced strategy adds one continuous auxiliary variable $z \ge 0$, the
worst assigned rank (Section~\ref{sec:obj}).

\subsection{Hard constraints}

These constraints are identical for every strategy; they are implemented once,
in the class \code{AssignmentModel} in \code{app/engine/solvers.py}.
\begin{align}
  \sum_{p \in P} x_{sp} &= 1 && \forall s \in S
    \label{eq:one}\\
  \sum_{s \in S} x_{sp} &\ge l_p && \forall p \in P
    \label{eq:min}\\
  \sum_{s \in S} x_{sp} &\le u_p \, y_p && \forall p \in P
    \label{eq:max}\\
  \sum_{s \in S} x_{sp} &\ge y_p && \forall p \in P
    \label{eq:runs}\\
  x_{sp} &= 1 && \forall (s,p) \in F
    \label{eq:fixed}\\
  \sum_{s \in S \,:\, \lambda_{sk} \ge L_{pk}} x_{sp} &\ge y_p
    && \forall p \in P,\ k \in M_p
    \label{eq:mand}\\
  \sum_{(s,p) \in A} x_{sp} &\le n - 1 && \forall A \in \mathcal{A}
    \label{eq:nogood}\\
  x_{sp},\, y_p &\in \{0,1\} && \forall s \in S,\ p \in P \notag
\end{align}

\begin{description}[leftmargin=0pt,font=\normalfont\bfseries]
  \item[\eqref{eq:one} One project per student.] Every student receives
        exactly one project.
  \item[\eqref{eq:min}--\eqref{eq:max} Capacity.] Each project receives
        between $l_p$ and $u_p$ students. The minimum applies to every project,
        consistent with input validation.
  \item[\eqref{eq:max}--\eqref{eq:runs} Linking $y_p$ to $x_{sp}$.] Together
        these force $y_p = 1$ exactly when project $p$ has at least one
        student. Constraint~\eqref{eq:max} is the big-$M$ linking constraint
        of the fixed-charge problem~\cite[Ch.~12.3]{ior} with the tightest
        valid $M = u_p$; constraint~\eqref{eq:runs} is a contingent decision
        (``$y_p$ only if some $x_{sp}$'').
  \item[\eqref{eq:fixed} Fixed assignments.] Administrative locks hold in
        every solution and every alternative.
  \item[\eqref{eq:mand} Mandatory skills.] A running project must include
        at least one student who meets each mandatory level. This is a
        set-covering constraint~\cite[Ch.~12]{ior}: the qualified students
        ``cover'' the requirement. Multiplying by $y_p$ means that a project
        with $l_p = 0$ may stay empty instead of requiring an expert. We
        interpret ``mandatory'' at team level (at least one qualified member);
        the stricter reading, every member qualified, would instead set
        $x_{sp} = 0$ for every unqualified student.
  \item[\eqref{eq:nogood} No-good cuts.] If every student kept the project
        they had in an earlier solution $A$, the left-hand side would equal
        $n$. Requiring at most $n - 1$ forces at least one student to move, so
        each new solution is distinct from all earlier ones. The constraint
        has the same structure as subtour elimination in the travelling
        salesman problem: it removes unwanted integer solutions and leaves all
        others feasible.
\end{description}

\subsection{Objectives: three strategies}
\label{sec:obj}

Each strategy is a subclass of \code{AssignmentStrategy} (Strategy pattern)
and differs only in its objective.

\paragraph{Strategy 1: preference-oriented (\code{PREFERENCE}).}
Minimise the total assigned rank; skill fit only breaks ties:
\begin{equation}
  \max \; Z_{\text{pref}} = \sum_{s \in S} \sum_{p \in P}
    \bigl(-r_{sp} + \varepsilon\, \sigma_{sp}\bigr)\, x_{sp}.
  \label{eq:zpref}
\end{equation}
Total rank changes in integer steps, while the tie-break term is at most
$\varepsilon n = 0.01$ for the demo cohort, so it can never outweigh one rank
step as long as $n < 1/\varepsilon = 1000$.

\paragraph{Strategy 2: skill-fit-oriented (\code{SKILL\_FIT}).}
Maximise total skill fit; preferences only break ties:
\begin{equation}
  \max \; Z_{\text{skill}} = \sum_{s \in S} \sum_{p \in P}
    \bigl(\sigma_{sp} + \varepsilon\, u_{sp}\bigr)\, x_{sp}.
  \label{eq:zskill}
\end{equation}

\paragraph{Strategy 3: balanced trade-off (\code{BALANCED}).}
A two-stage (lexicographic) objective. Stage~1 is a minimax problem, like the
districting problem in which the largest deviation is made as small as
possible: it finds the smallest achievable worst rank $z^*$ among students
without a fixed assignment,
\begin{equation}
  z^* = \min \; z \quad \text{s.t.} \quad
  z \ge r_{sp}\, x_{sp} \;\; \forall s \in S \setminus S_F,\ p \in P,
  \quad \eqref{eq:one}\text{--}\eqref{eq:nogood}.
  \label{eq:stage1}
\end{equation}
Fixed students are excluded because their rank is not a choice the model can
improve. Stage~2 keeps that guarantee and maximises an equal blend of
preference score and skill fit:
\begin{equation}
  \max \; Z_{\text{bal}} = \sum_{s \in S} \sum_{p \in P}
    \Bigl(\tfrac12\, u_{sp} + \tfrac12\, \sigma_{sp}\Bigr) x_{sp}
  \quad \text{s.t.} \quad
  x_{sp} = 0 \;\; \forall s \notin S_F,\ p : r_{sp} > z^*,
  \quad \eqref{eq:one}\text{--}\eqref{eq:nogood}.
  \label{eq:stage2}
\end{equation}

\subsection{Generating distinct alternatives}

\code{AssignmentEngine.generate\_alternatives} produces up to $K$ alternatives
(default $K = 3$):
\begin{enumerate}
  \item Alternative 1 uses the strategy the user selected; alternatives 2 and
        3 use the two other strategies, so each alternative favours a
        different criterion.
  \item Each solve includes a no-good cut~\eqref{eq:nogood} for every earlier
        solution, so all alternatives are pairwise distinct.
  \item The constraints are the same for every strategy, so if one solve is
        infeasible, no further distinct assignment exists. Generation then
        stops and reports the number found (\code{distinct\_alternatives}).
        It never fabricates duplicates.
  \item If the first solve is infeasible, the run is reported as
        \code{INFEASIBLE} with no solutions.
\end{enumerate}

\section{Evaluation indicators}
\label{sec:ind}

Let $p(s)$ be the project assigned to student $s$ in a solution, and
$\tilde r_s = r_{s\,p(s)}$ its rank (with $m+1$ for an unranked project). All
indicators are computed from the solution itself, independently of the solver.

\subsection{Feasibility}
The solution is re-checked against every hard constraint, so that a solver or
data error cannot pass unnoticed:
\begin{itemize}
  \item completeness: $\sum_{s,p} x_{sp} = n$, each student exactly once;
  \item capacity compliance: $l_p \le \sum_s x_{sp} \le u_p$ for every $p$;
  \item fixed-assignment adherence: $x_{sp} = 1$ for every $(s,p) \in F$;
  \item mandatory coverage: every running project has at least one qualified
        student for each $k \in M_p$.
\end{itemize}
Any violation is listed by name and marks the solution as infeasible.

\subsection{Efficiency}
\begin{itemize}
  \item \textbf{Mean skill fit}, our main efficiency measure:
        $\bar\sigma = \frac1n \sum_{s \in S} \sigma_{s\,p(s)} \in [0,1]$.
        It answers ``how well, on average, do students meet the levels their
        projects ask for?''. 1 means every requirement is fully met.
  \item \textbf{Unmet desired requirements}:
        $U_{\text{soft}} = \bigl|\{(s,k) : k \in K_{p(s)} \setminus M_{p(s)},\
        \lambda_{sk} < L_{p(s)k}\}\bigr|$, the number of desired (non-mandatory)
        skill gaps, which shows \emph{where} skill fit is lost.
  \item \textbf{Mean preference rank}:
        $\bar r = \frac1n \sum_{s \in S} \tilde r_s$, where lower is better.
\end{itemize}

\subsection{Fairness}
Aggregate performance is not enough, so we report how outcomes are
distributed:
\begin{enumerate}
  \item \textbf{Rank distribution}: $c_k = |\{s : \tilde r_s = k\}|$ for
        $k = 1,\dots,m$, plus $c_{\text{unranked}}$.
  \item \textbf{Students outside their preferences}:
        $U = |\{s : p(s) \text{ not ranked by } s\}|$.
  \item \textbf{Worst individual outcome}: $r_{\max} = \max_{s \in S} \tilde r_s$.
  \item \textbf{Our proposed indicator: the Gini coefficient of ranks},
        \begin{equation}
          G = \frac{\sum_{i \in S} \sum_{j \in S} |\tilde r_i - \tilde r_j|}
                   {2\, n \sum_{i \in S} \tilde r_i} \in [0, 1).
          \label{eq:gini}
        \end{equation}
        \emph{Justification.} $\bar r$ and $r_{\max}$ do not show whether
        disappointment is spread evenly or concentrated on a few students.
        $G$ does: it is 0 when every student receives the same rank, and
        grows as outcomes become more unequal. For ranks $(1,1,1,4)$, the mean
        rank is 1.75 and $G = 0.32$; for $(2,2,2,2)$ the mean rank is worse,
        2.00, but $G = 0$. Showing both lets the decision-maker see the
        efficiency--equity trade-off directly.
        \emph{Interpretation.} Because a low rank is good, a high $G$ means
        unequal \emph{disappointment}, not unequal wealth. $G$ measures
        spread, not level: ranks $(3,3,3,3)$ also give $G = 0$, so it is
        always reported together with $\bar r$ and $r_{\max}$.
        The rank variance $\frac1n \sum_s (\tilde r_s - \bar r)^2$ is
        reported as a supporting measure.
\end{enumerate}

\subsection{Capacity utilisation}
\begin{equation*}
  \text{Util}_p = \frac{\sum_{s} x_{sp}}{u_p} \times 100\%,
  \qquad
  \text{Util} = \frac{n}{\sum_{p} u_p} \times 100\%.
\end{equation*}
Projects at their minimum or maximum are flagged, since they leave no slack
when the situation changes.

\section{Trade-off analysis and solution comparison}

\subsection{Comparative deltas}
For a baseline solution $A$ and an alternative $B$, each indicator is reported
as a difference $B - A$:
\begin{equation*}
  \Delta\bar\sigma = \bar\sigma_B - \bar\sigma_A, \quad
  \Delta\bar r = \bar r_B - \bar r_A, \quad
  \Delta r_{\max} = r_{\max,B} - r_{\max,A}, \quad
  \Delta G = G_B - G_A.
\end{equation*}
A positive $\Delta\bar\sigma$ is an improvement; positive $\Delta\bar r$,
$\Delta r_{\max}$ and $\Delta G$ are deteriorations.

\subsection{Who is affected}
\begin{itemize}
  \item Changed assignments: $C = \{(s, p_A(s), p_B(s)) : p_A(s) \neq p_B(s)\}$;
        $|C| \ge 1$ for distinct alternatives.
  \item Worsened students: $\tilde r^B_s > \tilde r^A_s$; improved students:
        $\tilde r^B_s < \tilde r^A_s$.
\end{itemize}
This answers the brief's questions: what improves, what becomes worse, and who
is affected.

\subsection{Dominance}
Alternative $B$ is \emph{dominated} by $A$ if $A$ is at least as good on mean
rank and mean skill fit and strictly better on at least one
($\bar r_A \le \bar r_B$, $\bar\sigma_A \ge \bar\sigma_B$, not both equal).
A dominated alternative offers no real trade-off. Each strategy's optimum is
Pareto-optimal for its own criterion, but no-good cuts only guarantee
distinctness, so a later alternative can be dominated. Dominated alternatives
are therefore flagged: the engine adds a warning and the numbers of the
dominating alternatives (\code{diagnostics["dominated\_by"]}) to the
solution, and the comparison summary states that there is no real trade-off.

\subsection{Illustrative results}

\begin{table}[!htbp]
\centering
\caption{Standard demo cohort: 10 students, 4 projects with capacity 1--3
         (44 binary variables, each solve $< 0.1$\,s).}
\label{tab:std}
\begin{tabular}{lcccccccc}
\toprule
Alternative & $\bar r$ & $r_{\max}$ & $c_1$ & $U$ & $G$ & $\bar\sigma$ & $U_{\text{soft}}$ & Moved\\
\midrule
1 \code{PREFERENCE} & 1.00 & 1 & 10 & 0 & 0.000 & 0.983 & 1 & --\\
2 \code{SKILL\_FIT}  & 1.10 & 2 &  9 & 0 & 0.082 & 1.000 & 0 & 1 vs.\ alt.~1\\
3 \code{BALANCED}   & 1.10 & 2 &  9 & 0 & 0.082 & 0.983 & 1 & 2 vs.\ alt.~2\\
\bottomrule
\end{tabular}
\end{table}

Table~\ref{tab:std} shows a clean trade-off between alternatives 1 and 2:
moving a single student (S06) to their second choice raises mean skill fit
from 0.983 to 1.000 and removes the only unmet desired requirement. Alternative
3 is dominated by both: it has a worse mean rank than alternative 1 at the same
skill fit, and the same mean rank as alternative 2 at a lower skill fit. The
engine therefore attaches the warning ``Dominated by alternatives 1 and 2''. Because alternative 1 already gives everyone a
first choice, the balanced strategy cannot improve the worst rank in this
cohort; it only becomes a different alternative because of the no-good cuts.

\begin{table}[!htbp]
\centering
\caption{Fixed-assignment scenario: S02 is locked to P2, their third choice.}
\label{tab:fixed}
\begin{tabular}{lcccccc}
\toprule
Alternative & $\bar r$ & $r_{\max}$ & $c_1/c_2/c_3$ & $G$ & $\bar\sigma$ & $U_{\text{soft}}$\\
\midrule
1 \code{PREFERENCE} & 1.20 & 3 & 9/0/1 & 0.150 & 0.967 & 2\\
2 \code{SKILL\_FIT}  & 1.30 & 3 & 8/1/1 & 0.192 & 0.983 & 1\\
3 \code{BALANCED}   & 1.30 & 3 & 8/1/1 & 0.192 & 0.967 & 2\\
\bottomrule
\end{tabular}
\end{table}

In Table~\ref{tab:fixed} the lock holds in every alternative, and its cost is
visible: compared with Table~\ref{tab:std}, mean rank rises from 1.00 to 1.20
and the worst rank is 3, caused entirely by the fixed student.

The demo cohort is too easy to show the balanced strategy at work. In a
4-student, 4-project example with one seat per project (preferences
S1, S2: D,C,A,B; S3: A,C,B,D; S4: C,A,D,B), \code{PREFERENCE} reaches the
lowest total rank, 7 ($\bar r = 1.75$), by giving S2 their fourth choice
($r_{\max} = 4$, $G = 0.32$). \code{BALANCED} accepts a higher total of 8
($\bar r = 2.00$) to guarantee everyone their third choice or better
($r_{\max} = 3$, $G = 0.19$). This is precisely the trade-off the brief
describes: protecting the worst-off student reduces average quality.

\section{Transparent recommendation rules}

Recommendations are generated in \code{app/engine/recommendations.py} for each
policy priority separately. Each recommendation lists the criteria used, the
calculated evidence (indicator values) and the advice, so that calculated
facts are kept apart from normative judgement.

\begin{table}[!htbp]
\centering
\small
\caption{Recommendation rules. Criteria are applied lexicographically: the
         second criterion only decides between alternatives tied on the
         first.}
\begin{tabularx}{\textwidth}{>{\raggedright\arraybackslash}p{3.1cm}LL}
\toprule
Policy priority & Selection criterion (lexicographic) & Advice and trade-off warning\\
\midrule
1. Student preference maximisation
  & minimise $U$, then $\bar r$, then $r_{\max}$
  & Adopt when satisfaction and few appeals matter most; may give projects
    weaker skill coverage.\\
2. Project skill quality and delivery
  & maximise $\bar\sigma$, then minimise $U_{\text{soft}}$
  & Adopt when partners demand strong technical capability; some students
    receive lower-ranked projects.\\
3. Balanced and egalitarian
  & minimise $G$, then rank variance, then $\bar r$
  & Adopt to avoid concentrating disappointment on a few students; may
    lower average satisfaction.\\
\bottomrule
\end{tabularx}
\end{table}

When several alternatives tie on all criteria, the one with the lowest
alternative number is recommended. The same alternative can be recommended
under more than one priority; this is shown rather than hidden, because it
tells the decision-maker that the priorities do not conflict for this cohort.

\section{Assumptions, boundary conditions and edge cases}

\paragraph{Assumptions.}
\begin{enumerate}
  \item \textbf{Unranked projects} are allowed but cost rank $m + 1$, worse
        than any ranked project, and are counted in $U$.
  \item \textbf{Ranks are evenly spaced}: the step from first to second
        choice weighs the same as from second to third.
  \item \textbf{Skill levels} are taken at face value; missing skills count
        as level 0. A project with no requirements has fit 1.
  \item \textbf{Mandatory skills} are met at team level (at least one
        qualified member), pending the instructor's confirmation.
  \item \textbf{Ties} between students with identical preferences are broken
        by the secondary goal ($\varepsilon$ term) and then by the solver; the
        result is a valid optimum, not a judgement about which student
        deserves the seat.
\end{enumerate}

\paragraph{Handling of edge cases.}
Every case in the brief has a defined outcome; none may end in an unexplained
crash. Input validation (\code{app/domain/validation.py}) runs before the
solver.

\begin{table}[!htbp]
\centering
\small
\caption{Edge cases and the system's response.}
\label{tab:edge}
\begin{tabularx}{\textwidth}{>{\raggedright\arraybackslash}p{5.4cm}>{\raggedright\arraybackslash}p{2.3cm}L}
\toprule
Case & Outcome & Behaviour\\
\midrule
Empty cohort or catalogue; duplicate or unknown identifiers; invalid
  capacities ($u_p \le 0$, $l_p < 0$, $l_p > u_p$)
  & Reject & Named validation error; the solver is not run.\\
Contradictory fixed assignments (one student fixed twice)
  & Reject & Named validation error.\\
Incomplete, duplicate or tied preferences; no preferences
  & Warn & Run proceeds; unranked projects cost $m+1$.\\
Uncatalogued skill in a profile or requirement
  & Warn & Run proceeds; the skill counts only where defined.\\
Missing skills; project without requirements
  & Accept & Fit 0 or 1 respectively (equation~\eqref{eq:fit}).\\
$\sum_p u_p < n$; $\sum_p l_p > n$; more fixed students than $u_p$;
  a mandatory skill no student has
  & Infeasible & Proven before solving, with the reason stated.\\
Infeasibility arising only from combinations (e.g.\ two projects each need
  the single qualified student)
  & Infeasible & Proven by the solver; status \code{INFEASIBLE}, no
    solutions.\\
Only one (or fewer than $K$) feasible solutions
  & Partial & The solutions found are returned; the stop reason states that
    no further distinct solution exists.\\
Duplicate alternatives
  & Prevented & Impossible by construction (no-good cuts).\\
Equal aggregate scores
  & Accept & Tie-break term; in recommendations the lowest alternative number wins.\\
Solver time limit (15\,s) or database failure
  & Fail & Run status \code{NOT\_SOLVED}/\code{FAILED} with diagnostics.\\
\bottomrule
\end{tabularx}
\end{table}

\paragraph{Feasibility boundaries.}
The problem is provably infeasible if $\sum_p u_p < n$, if
$\sum_p l_p > n$, if any project has more fixed students than $u_p$, or if a
mandatory level is held by no student in the cohort. Validation detects these
directly. Validation is deliberately conservative: it treats the last case as
infeasible even when the project has $l_p = 0$ and could stay empty under
constraint~\eqref{eq:mand}. All other infeasibility, which arises from
combinations of constraints, is proven by the solver.

\section{Limitations}

\begin{enumerate}
  \item \textbf{Computational complexity.} Integer programming is NP-hard in
        general, and the number of binary variables grows as $n \cdot m$. At
        realistic cohort sizes the model remains small, but we have only
        measured solve times on the demo cohorts.
  \item \textbf{Linear preferences.} Evenly spaced ranks assume that a
        student is as disappointed moving from first to second choice as
        from second to third, which need not hold.
  \item \textbf{Individual, not team, skill fit.} $\sigma_{sp}$ is measured
        per student. The model does not capture complementary skills, team
        diversity or personal compatibility.
  \item \textbf{Weights and tie-breaks are judgements.} The $\tfrac12/\tfrac12$
        blend in \code{BALANCED} is a choice, not a fact. In
        \code{SKILL\_FIT} the tie-break term can outweigh very small
        differences in total skill fit (below $\varepsilon n$), and stage~2 of
        \code{BALANCED} has no tie-break, so the solver chooses between
        equally good assignments. We have not run a sensitivity analysis on
        these parameters~\cite[Ch.~7]{ior}.
  \item \textbf{Alternatives are distinct, not necessarily efficient.} No-good
        cuts guarantee that alternatives differ, not that they are
        Pareto-optimal. Dominated alternatives are flagged
        (Table~\ref{tab:std}) but still generated.
  \item \textbf{Static, self-reported data.} The model assumes the input data
        are correct and fixed. Self-reported skill levels may be inaccurate,
        and mid-semester changes require a new run.
\end{enumerate}

\begin{thebibliography}{9}
\bibitem{rezaeinia2022}
N.~Rezaeinia, J.~C. G\'oez and M.~Guajardo.
\newblock Efficiency and fairness criteria in the assignment of students to
  projects.
\newblock \emph{Annals of Operations Research}, 319(2):1717--1735, 2022.
\newblock \url{https://doi.org/10.1007/s10479-021-04001-7}

\bibitem{ior}
F.~S. Hillier and G.~J. Lieberman.
\newblock \emph{Introduction to Operations Research}.
\newblock McGraw-Hill. Chapters 7 (sensitivity analysis), 9.3 (the
  assignment problem) and 12 (integer programming), as covered in TECH4
  Mathematics for Decisions, NHH, autumn 2025.

\bibitem{pyomo}
M.~L. Bynum, G.~A. Hackebeil, W.~E. Hart, C.~D. Laird, B.~L. Nicholson,
  J.~D. Siirola, J.-P. Watson and D.~L. Woodruff.
\newblock \emph{Pyomo: Optimization Modeling in Python}, 3rd edition.
\newblock Springer, 2021.

\bibitem{highs}
Q.~Huangfu and J.~A.~J. Hall.
\newblock Parallelizing the dual revised simplex method.
\newblock \emph{Mathematical Programming Computation}, 10(1):119--142, 2018.
\end{thebibliography}

\end{document}
